\documentclass[twocolumn,11pt]{article}

\usepackage{natbib}
\usepackage{amsmath}
\usepackage{geometry}
\usepackage{graphicx}
\usepackage{hyperref}
\usepackage{times}
\usepackage{enumitem}
\usepackage[switch]{lineno}

\geometry{a4paper, margin=0.75in}
\urlstyle{same}

\begin{document}

\title{Consistent Anime Character Sheet Generation from 1--3 Reference Images\\Midway Executive Summary}
\author{
Elaine Zhang (\texttt{<andrew-email>}) \and
Sijie Liu (\texttt{<andrew-email>}) \and
Zihui Jin (\texttt{<andrew-email>})\\
10-423/623 Generative AI Course Project
}
\date{April 13, 2026}
\maketitle
\small

\section{Introduction}
We study a practical failure mode of modern text-to-image systems: they can generate high-quality anime portraits, but they do not reliably produce a \emph{character sheet}---a multi-panel reference sheet with consistent identity across expressions (and optionally views). Our project targets \textbf{expression-conditioned identity-preserving generation}: given 1--3 reference images of a new OC and a target expression, the model should synthesize the same character with the requested expression. We propose a reference-guided pipeline that fine-tunes an IP-Adapter backbone using paired expression data, plus additional control to reduce identity drift and copy failures. At midway, we have completed the expression-sheet dataset pipeline, implemented baseline inference and evaluation, and prepared scripts for prompt-only, Textual Inversion, LoRA (global expression control), and IP-Adapter baselines. We expect our reference-guided fine-tuned model to improve identity consistency and expression accuracy over prompt-only and global prompt-conditioned baselines.

\section{Dataset and Task}
\textbf{Dataset.} We collect expression sheets from SafeBooru (SFW Danbooru) using tags such as \texttt{expression\_sheet}, \texttt{multiple\_expressions}, and \texttt{expressions}. Images are filtered by size and mode, then face crops are extracted using an anime-face detector. Crops from the same sheet are treated as the same character. We create paired examples from the same sheet and split by \texttt{sheet\_id} to prevent leakage. At midpoint, we have a cleaned dataset and a held-out split ready for baseline experiments.

\textbf{Task.} Stage-1 task is expression-controlled identity generation:
\[
(x_{\mathrm{ref}}^{1:K}, e_{\mathrm{target}})\rightarrow \hat{x}_{\mathrm{expr}},
\]
where $x_{\mathrm{ref}}^{1:K}$ are 1--3 reference images of one OC and $e_{\mathrm{target}} \in \{\text{neutral, happy, sad, angry, surprised, crying}\}$.

\textbf{Labels.} We use CLIP zero-shot labeling to assign an expression label to each crop; these labels are used for supervised baselines and evaluation.

\textbf{Evaluation.} We evaluate identity preservation (ArcFace cosine similarity), color consistency (Lab palette distance), image quality (FID), expression accuracy (CLIP classifier), copy score (SSIM-based), and sheet-level correctness (coverage and panel accuracy). These metrics are implemented in our evaluation scripts and used across baselines.

\section{Related Work}
Latent diffusion models are the backbone for controllable generation \citep{rombach2022latent}. Personalization methods such as Textual Inversion and LoRA provide few-shot adaptation but require per-character training \citep{gal2022textual,hu2022lora}. IP-Adapter injects reference-image features into diffusion models and is a natural reference-guided baseline \citep{ye2023ipadapter}. ControlNet provides explicit structural conditioning, useful for multi-view extensions \citep{zhang2023controlnet}. Recent identity-preserving approaches motivate multi-reference conditioning and stronger identity constraints \citep{li2024photomaker,guo2024pulid,ma2024characteradapter}. Story-level consistency methods emphasize persistent memory across images \citep{zhou2024storydiffusion}. Copy-failure analyses highlight the need to penalize trivial reference copying \citep{xu2025withanyone}. These works inform our pipeline design for anime character sheets.

\section{Approach}
\textbf{Baselines.} We compare three baseline families on the same expression task:
\begin{itemize}[leftmargin=*,nosep]
    \item \textbf{Prompt-only baseline} (teammate result pending): structured prompts vs.\ naive prompts on the same SD1.5 backbone.
    \item \textbf{Textual Inversion baseline}: per-character token embedding learned from a few samples.
    \item \textbf{Global LoRA baseline}: a single LoRA adapter trained on all labeled expression pairs to learn prompt-based expression control that transfers to new characters.
    \item \textbf{Vanilla IP-Adapter baseline}: zero-shot reference-guided generation using SD1.5 + IP-Adapter.
\end{itemize}

\textbf{Main method (to implement).} We will fine-tune IP-Adapter on paired expression data. The model uses reference images for identity and an emotion prompt for control. We will explore multi-reference fusion and anti-copy regularization in training, then assemble panels into a final OC sheet. This preserves the core reference-guided design while improving expression correctness.

\section{Experiments}
\textbf{Baseline results.} We have run the vanilla IP-Adapter baseline on held-out pairs and will report identity, palette, FID, and expression accuracy. Prompt-only baseline results will be added by a teammate.

\textbf{Planned experiments.} We will compare:
\begin{itemize}[leftmargin=*,nosep]
    \item prompt-only vs.\ global LoRA vs.\ vanilla IP-Adapter vs.\ IP-Adapter fine-tune,
    \item single-reference vs.\ multi-reference conditioning (ablation),
    \item with vs.\ without anti-copy regularization.
\end{itemize}

\textbf{Skeleton table.} Table~\ref{tab:results} will contain quantitative metrics for all baselines and our method. Prompt-only results are intentionally left blank for teammate completion.

\begin{table}[h]
    \centering
    \small
    \begin{tabular}{|l|c|c|c|c|c|}
        \hline
        \textbf{Method} & \textbf{FID} $\downarrow$ & \textbf{ArcFace} $\uparrow$ & \textbf{Palette} $\downarrow$ & \textbf{Expr. Acc.} $\uparrow$ & \textbf{Copy} $\downarrow$ \\
        \hline
        Prompt-only & -- & -- & -- & -- & -- \\
        Textual Inversion & -- & -- & -- & -- & -- \\
        Global LoRA & -- & -- & -- & -- & -- \\
        IP-Adapter (zero-shot) & -- & -- & -- & -- & -- \\
        IP-Adapter (fine-tuned) & -- & -- & -- & -- & -- \\
        \hline
    \end{tabular}
    \caption{Planned quantitative results. Prompt-only row intentionally left blank for teammate results.}
    \label{tab:results}
\end{table}

\textbf{Qualitative examples.} Figure~\ref{fig:qual} will show reference + generated panels for each baseline.

\begin{figure}[h]
    \centering
    \fbox{\parbox{0.92\linewidth}{\centering Placeholder: qualitative panels for each baseline}}
    \caption{Qualitative comparison (reference + generated expressions).}
    \label{fig:qual}
\end{figure}

\section{Plan, Human Work, and AI Use}
\textbf{Remaining timeline.}
\begin{itemize}[leftmargin=*,nosep]
    \item \textbf{Apr 13--Apr 20:} Run prompt-only baseline (Member 2), finalize Textual Inversion + LoRA baselines (Member 1/3).
    \item \textbf{Apr 20--Apr 27:} IP-Adapter fine-tuning run and ablations (Member 3).
    \item \textbf{Apr 27--May 3:} Evaluate baselines and main method, finalize figures/tables (All).
    \item \textbf{May 3--May 7:} Write final report and polish.
\end{itemize}

\textbf{Human work.} Dataset curation, face-crop filtering, and baseline experiment design were performed by the team. Qualitative checks and metric selection were decided by the group.

\textbf{AI use.} We used AI assistance to draft code templates (e.g., CLIP labeling, LoRA training scaffolds) and to rewrite sections for clarity. We verified the generated code by running it on sample data, checking file outputs, and comparing metrics across baselines. No AI was used to generate or manipulate dataset images.

\textbf{Future AI use.} We plan to use AI assistance for documentation polishing and for small script refactors, while keeping training, evaluation, and analysis under direct human control.

\section{Thought-Experiment on Compute}
\textbf{Actual compute so far.} We have run baseline inference on a held-out set on a single GPU (type: \texttt{<GPU type>}) for approximately \texttt{<X>} GPU-hours. We will update the final summary with the exact compute and estimated dollar cost using the course price sheet.

\textbf{If given \$1{,}000 in GPU credits.} We would (1) scale up IP-Adapter fine-tuning with longer schedules and larger batch sizes, (2) run a robust multi-reference ablation study, and (3) attempt the multi-view extension with ControlNet line art to complete full OC sheets.

\clearpage
\bibliographystyle{plainnat}
\bibliography{references}

\end{document}
